\documentclass[11pt,a4paper]{article}
\usepackage{fontspec,amsmath,amssymb,amsthm,microtype}
\setmainfont{lmroman10-regular.otf}[BoldFont=lmroman10-bold.otf,ItalicFont=lmroman10-italic.otf,BoldItalicFont=lmroman10-bolditalic.otf]
\usepackage{xcolor}
\usepackage[most]{tcolorbox}
\usepackage{mathrsfs}
\providecommand{\lt}{<}
\providecommand{\gt}{>}
\setmonofont{DejaVuSansMono.ttf}[BoldFont=DejaVuSansMono-Bold.ttf,ItalicFont=DejaVuSansMono-Oblique.ttf,BoldItalicFont=DejaVuSansMono-BoldOblique.ttf,Scale=0.86]
\newfontfamily\leanfallback{STIXTwoMath-Regular.otf}[Scale=0.86]
\definecolor{leanrule}{gray}{0.62}
\definecolor{leanback}{gray}{0.962}
\definecolor{leantitle}{gray}{0.22}
\definecolor{leanmuted}{gray}{0.38}
\newtcolorbox{leanbox}[1]{enhanced,breakable,colback=leanback,colframe=leanrule,
  boxrule=0.4pt,arc=0.6mm,left=2.4mm,right=2.4mm,top=1.8mm,bottom=1.8mm,
  before skip=0.8em,after skip=1em,fonttitle=\small\bfseries,coltitle=leantitle,
  colbacktitle=leanback,titlerule=0.3pt,title={#1}}
\newlength{\leancharwidth}
\newenvironment{leancode}{\par\smallskip\begingroup\ttfamily\footnotesize\raggedright
  \settowidth{\leancharwidth}{M}\setlength{\parskip}{0pt}\setlength{\parindent}{0pt}
  \frenchspacing}{\par\endgroup\smallskip}
\newcommand{\leanline}[2]{\hangindent=\dimexpr#1\leancharwidth+4\leancharwidth\relax
  \hangafter=1\noindent\hspace*{#1\leancharwidth}#2\par}
\newcommand{\leansource}[1]{\par{\footnotesize\color{leanmuted}#1}\par}
\newcommand{\leanlabel}[1]{\par\smallskip\noindent{\small\bfseries #1}\par}
\usepackage[margin=25mm]{geometry}
\usepackage[unicode,hidelinks]{hyperref}
\hypersetup{pdfauthor={LangevinLean Project}}
\setlength{\emergencystretch}{3em}
\setlength{\parskip}{0.4em}
\setlength{\parindent}{0pt}
\allowdisplaybreaks[2]
\raggedbottom
\title{\texorpdfstring{The actual reversible strongly continuous contraction semigroup}{The actual reversible strongly continuous contraction semigroup}}
\author{LangevinLean Project}
\date{}
\begin{document}\maketitle
Draft mathematical exposition; these notes may contain mistakes.


All operator properties required by the corrector construction are assembled for the actual Gibbs Langevin diffusion. The semigroup law, self-adjointness, contraction, identity at zero and strong continuity are proved from the actual process.


\section{\texorpdfstring{Assembling the actual operator properties}{Assembling the actual operator properties}}


\textbf{Claim.} The canonical transition family is a reversible strongly continuous contraction semigroup on complex \(L^2(\pi)\).



\textbf{Proof.} The actual conditional Markov theorem and compression calculation prove the real composition law. Scalar restriction identifies the real and complex operators, so that law transfers to complex observables. Reversibility follows from actual reversible MALA approximants and their strong limit; contraction, the time-zero identity and strong continuity follow from actual stationary evaluation. Each defining property is therefore a theorem about the same canonical Brownian operator family.

\begin{leanbox}{Lean: \texttt{op\_\allowbreak{}add\_\allowbreak{}real}, \texttt{instSemigroupInstanceReal}, \texttt{op\_\allowbreak{}add} and 2 more}
\leansource{\texttt{op\_\allowbreak{}add\_\allowbreak{}real} \quad WassersteinMixing/\allowbreak{}SemigroupInstance.lean, lines 26–31}
\begin{leancode}
\leanline{0}{theorem\ op\_add\_real\ (P\ :\ Parameters)\ (V\ :\ E\ P.dimension\ →\ ℝ)}
\leanline{4}{(hV\ :\ PotentialAssumptions\ P\ V)\ (s\ t\ :\ ℝ)\ (hs\ :\ 0\ ≤\ s)\ (ht\ :\ 0\ ≤\ t)\ :}
\leanline{4}{op\ ({\leanfallback 𝕜}\ :=\ ℝ)\ (F\ :=\ F)\ P\ V\ hV\ (s\ +\ t)\ =\ op\ P\ V\ hV\ s\ ∘L\ op\ P\ V\ hV\ t}
\end{leancode}
\leansource{\texttt{instSemigroupInstanceReal} \quad WassersteinMixing/\allowbreak{}SemigroupInstance.lean, lines 37–38}
\begin{leancode}
\leanline{0}{local\ instance\ instSemigroupInstanceReal\ :\ InnerProductSpace\ ℝ\ F\ :=}
\leanline{2}{InnerProductSpace.rclikeToReal\ ℂ\ F}
\end{leancode}
\leansource{\texttt{op\_\allowbreak{}add} \quad WassersteinMixing/\allowbreak{}SemigroupInstance.lean, lines 42–51}
\begin{leancode}
\leanline{0}{theorem\ op\_add\ (P\ :\ Parameters)\ (V\ :\ E\ P.dimension\ →\ ℝ)}
\leanline{4}{(hV\ :\ PotentialAssumptions\ P\ V)\ (s\ t\ :\ ℝ)\ (hs\ :\ 0\ ≤\ s)\ (ht\ :\ 0\ ≤\ t)\ :}
\leanline{4}{op\ ({\leanfallback 𝕜}\ :=\ ℂ)\ (F\ :=\ F)\ P\ V\ hV\ (s\ +\ t)\ =\ op\ P\ V\ hV\ s\ *\ op\ P\ V\ hV\ t}
\end{leancode}
\leansource{\texttt{semigroup} \quad WassersteinMixing/\allowbreak{}SemigroupInstance.lean, lines 55–62}
\begin{leancode}
\leanline{0}{def\ semigroup\ (P\ :\ Parameters)\ (V\ :\ E\ P.dimension\ →\ ℝ)\ (hV\ :\ PotentialAssumptions\ P\ V)\ :}
\leanline{4}{ReversibleSemigroup\ (Lp\ F\ 2\ (ULA.gibbsMeasure\ V))\ where}
\leanline{2}{op\ :=\ op\ P\ V\ hV}
\leanline{2}{zero\ :=\ op\_zero\ P\ V\ hV}
\leanline{2}{add\ :=\ fun\ \{s\ t\}\ hs\ ht\ ↦\ op\_add\ P\ V\ hV\ s\ t\ hs\ ht}
\leanline{2}{selfAdjoint\ :=\ fun\ \{t\}\ ht\ ↦\ op\_isSelfAdjoint\ P\ V\ hV\ t\ ht}
\leanline{2}{norm\_le\ :=\ fun\ \{t\}\ \_\ f\ ↦\ op\_norm\_le\ P\ V\ hV\ t\ f}
\leanline{2}{continuous\ :=\ fun\ f\ ↦\ continuousOn\_op\ P\ V\ hV\ f}
\end{leancode}
\leansource{\texttt{semigroup\_\allowbreak{}op} \quad WassersteinMixing/\allowbreak{}SemigroupInstance.lean, lines 64–66}
\begin{leancode}
\leanline{0}{@[simp]\ theorem\ semigroup\_op\ (P\ :\ Parameters)\ (V\ :\ E\ P.dimension\ →\ ℝ)}
\leanline{4}{(hV\ :\ PotentialAssumptions\ P\ V)\ (t\ :\ ℝ)\ :}
\leanline{4}{(semigroup\ (F\ :=\ F)\ P\ V\ hV).op\ t\ =\ op\ ({\leanfallback 𝕜}\ :=\ ℂ)\ P\ V\ hV\ t}
\end{leancode}
\leanlabel{How the Lean proof follows the argument}
\texttt{op\_\allowbreak{}add\_\allowbreak{}real} instantiates the actual Markov composition theorem. The local real structure and \texttt{op\_\allowbreak{}restrict\allowbreak{}Scalars} transfer it in \texttt{op\_\allowbreak{}add}. \texttt{semigroup} fills every field of \texttt{Reversible\allowbreak{}Semigroup} with the proved canonical operator results, and \texttt{semigroup\_\allowbreak{}op} records its definitional identification.
\end{leanbox}


Earlier results used here: \cite[Section~3.26.3]{LangevinLean2026ULASC}.

\begin{thebibliography}{ULA-LSI}
\bibitem[ULA-SC]{LangevinLean2026ULASC}
LangevinLean Project. \emph{Unadjusted Langevin under Strong Convexity}. LL-01, 2026. Draft report, content revision d34ead06987c. \url{https://odakyildiz.github.io/LangevinLean/docs/pdf/UnadjustedLangevin/FullProof.pdf}.
\end{thebibliography}

\end{document}
